% abstract/05_annotation_en
% Budget: 2 pages. Mandatory by section 9 (д), and it goes on the last page.
%
% The mechanism, checked on a probe rather than assumed:
%   \bgfalse scoped around this fragment does NOTHING. \bgen is defined once at load
%   time by branching on \ifbg, which is exactly the property that makes it expandable
%   and safe inside \addcontentsline. Flipping the conditional afterwards does not
%   redefine the macro.
%   \renewcommand{\bgen}[2]{#2} inside a group DOES work.
%   \begin{otherlanguage}{english} switches babel's part: hyphenation and the
%   automatic names, so \figurename becomes Figure inside it and reverts outside.
% Both are used below. The \bgen call is deliberate: it keeps the mechanism exercised
% by every build, so a preamble change that breaks it fails here rather than at the
% printer.

\clearpage
\begin{otherlanguage}{english}
\begingroup
\renewcommand{\bgen}[2]{#2}

\section*{\bgen{АНОТАЦИЯ}{ANNOTATION}}

The dissertation studies the number and organisation of listening descriptors in a
web application library that serves HTTP/1.1, HTTP/2, HTTP/3 and WebSocket at the
same time.

Conventional servers bind the choice of protocol to the choice of listening socket.
The choice is recorded when the socket is configured, that is, before any byte has
arrived, so each additional protocol adds a listening unit. The dissertation shows
this is documented behaviour across five servers and three C++ libraries rather than
an assumption, and identifies the one system that departs from it together with what
that departure costs.

The work moves the decision to the moment after the connection is accepted. A single
octet separates a TLS record from a cleartext method token, so one descriptor can
carry every application-layer protocol above a given transport. The number of
listening descriptors then depends on the accept model the platform requires and on the
number of transports, and not on the number of protocols served. This is stated as a
proposition with explicit assumptions and proved, so that it can be refuted.

Three further results support it. QUIC connections are distributed across worker
threads by encoding the worker index in the connection identifier, which is portable
to platforms where eBPF does not exist. Values that have not yet arrived are
expressed as a type on the server and as a matching primitive in the browser, so
incompleteness is carried by the type system on both sides of the connection rather
than by a convention. A second input and output mechanism on Linux makes the choice
between readiness and completion an independent variable within one codebase, one
machine and one kernel.

The measurement methodology is built around the ways a server benchmark can be
silently wrong, since an incorrectly conducted measurement does not refuse to run: it
produces plausible numbers. Both arms of every comparison come from one binary
selected by a runtime flag, the environment is fingerprinted and a change in it stops
the campaign, and the rejection criteria are declared before any run.

One campaign of three designs, \R{campaign.x1.runs}, \R{campaign.transport.runs} and
\R{campaign.churn.runs} runs, was carried out on one unvirtualised machine under Windows, for
HTTP/1.1 in cleartext and over TLS. \R{campaign.x1.rejected}, \R{campaign.transport.rejected}
and \R{campaign.churn.rejected} runs respectively were refused by the admissibility rules, each
with its reason recorded, and the worst value each rule reached is reported alongside those
counts.

The first hypothesis is not rejected. Across five offered loads from 10\,000 to
70\,000 requests per second, the difference in the 99th percentile between the
classifying arm and the same binary without classification lies between
\R{h1detect.x1.rel-min}\% and \R{h1detect.x1.rel-max}\%, and the confidence interval on
the difference covers zero at every load. The resolution of the design is reported with
the result: the half-width of that interval is between
\R{h1detect.x1.resolution-min}\% and \R{h1detect.x1.resolution-max}\% of the median, so
the design can see a difference below the five percent threshold it declared in advance
at four of the five loads, and not at 25\,000 requests per second. A non-rejection without that figure is a statement about the
number of runs rather than about the server.

The scope of the result is narrower than the hypothesis. Classification costs something
once per connection and something once per request, and a design holding sixty-four
persistent connections that carry between two hundred thousand and one million four
hundred thousand requests amortises the per-connection cost below any resolution. What is
bounded there is the per-request path. The per-connection cost was measured instead by the
churn design of the same campaign, \R{campaign.churn.runs} runs under connection-arrival
load; what remains is to repeat it over a network path.

Counting the listening descriptors of a running process, from the kernel rather than
from the server, confirms the proposition and corrects it. The count does not change
when classification is switched off, which is the claim the work defends. It also does
not change with the number of worker threads under Windows, which refutes the earlier
form of the equation and yields the multiplier taken from the platform's accept model.

The three remaining hypotheses are not tested here, for three different reasons, each
stated with the measurement that would settle it.

\vspace{0.5em}
\noindent
\emph{Keywords:} protocol demultiplexing, listening sockets, HTTP/3, QUIC connection
identifier, C++ coroutines, io\_uring, epoll, server benchmarking methodology,
coordinated omission, deferred rendering.

\endgroup
\end{otherlanguage}
